\documentclass{article}
\usepackage[utf8]{inputenc}
\usepackage[T1]{fontenc}
\usepackage{listings}
\usepackage{graphicx}
\usepackage{color}
\usepackage[top=15mm, bottom=10mm, left=13mm, right=13mm]{geometry}
\pagenumbering{gobble}
\usepackage{textgreek}

% Packages for headers and footers
\usepackage{fancyhdr}
\pagestyle{fancy}

% Define your custom R style
\definecolor{codegreen}{rgb}{0,0.6,0}
\definecolor{codegray}{rgb}{0.5,0.5,0.5}
\definecolor{codepurple}{rgb}{0.58,0,0.82}
\definecolor{backcolour}{rgb}{0.95,0.95,0.92}

\lstdefinestyle{mystyle}{
    backgroundcolor=\color{backcolour},   
    commentstyle=\color{codegreen},
    keywordstyle=\color{magenta},
    numberstyle=\tiny\color{codegray},
    stringstyle=\color{codepurple},
    basicstyle=\ttfamily\footnotesize,
    breakatwhitespace=false,         
    breaklines=true,                 
    keepspaces=true,                 
    numbers=left,                    
    numbersep=5pt,                  
    showspaces=false,                
    showstringspaces=false,
    showtabs=false,                  
    tabsize=2,
    literate=%
        {á}{{\'a}}1
        {é}{{\'e}}1
        {í}{{\'i}}1
        {ó}{{\'o}}1
        {ú}{{\'u}}1
        {Á}{{\'A}}1
        {É}{{\'E}}1
        {Í}{{\'I}}1
        {Ó}{{\'O}}1
        {Ú}{{\'U}}1
        {à}{{\`a}}1
        {è}{{\`e}}1
        {ì}{{\`i}}1
        {ò}{{\`o}}1
        {ù}{{\`u}}1
        {À}{{\`A}}1
        {È}{{\`E}}1
        {Ì}{{\`I}}1
        {Ò}{{\`O}}1
        {Ù}{{\`U}}1
        {ã}{{\~a}}1
        {õ}{{\~o}}1
        {Ã}{{\~A}}1
        {â}{{\~A}}1
        {Õ}{{\~O}}1
        {ç}{{\c c}}1
        {Ç}{{\c C}}1
        {α}{{\textalpha}}1
        {β}{{\textbeta}}1
        {μ}{{\textmu}}1
        {σ}{{\textsigma}}1
}

\lstset{style=mystyle}

% Define your headers
\rhead{Miguel Freitas (102756)}
\lhead{Afonso Ribeiro (102763)} % left side empty
\chead{Luís Correia (103528)} % center empty



\begin{document}
\centering
\noindent{\Large{\textbf{Projeto Computacional de Probabilidade e 
Estatística – Exercício 8}}}
\vspace{3mm}

\begin{lstlisting}[language=R, xleftmargin=4em, xrightmargin=4em]
# Configurar a semente
set.seed(1208)

# Parâmetros da distribuição de Cauchy
α <- 2.8
β <- 2

# Tamanho da amostra
n <- 163

# Parâmetros da distribuição normal
μ <- 0.5
σ <- sqrt(1.4)

# Calcular as probabilidades dos quantis
q <- seq(1, n) / (n + 1)

# Gerar a amostra e ordenar
amostra_cauchy <- sort(rcauchy(n, location = α, scale = β))

# Quantis da distribuição de Cauchy
quantis_cauchy <- qcauchy(q, location = α, scale = β)

# Quantis da distribuição Normal
quantis_normal <- qnorm(q, mean = μ, sd = σ)

# Criar o gráfico
dev.new()
par(mar = c(5, 5, 4, 2))
plot(quantis_cauchy, amostra_cauchy, col = "blue", pch = 19,
     xlab = "Quantis de probabilidade",
     ylab = "Amostra ordenada",
     cex.lab = 1.4, cex.main = 1.3, cex.axis = 1.3,
     main = paste("Gráfico Q-Q de uma amostra de uma população com distribuição",
                  "de Cauchy para distribuições Normal e de Cauchy"))

points(quantis_normal, amostra_cauchy, col = "red", pch = 19)
abline(0, 1, col = "green")
legend("bottomright", legend = c("Cauchy", "Normal"), col = c("blue", "red"), pch = 19, cex = 1.2)
\end{lstlisting}

\vspace{0.4cm}

\begin{figure}[ht]
\centering
\includegraphics[width=\linewidth, height=\textheight, keepaspectratio]{ex8.jpeg} 
\label{fig:mylabel}
\end{figure}

\end{document}
